% ============================================================
% data/achievements.tex
% Defines all achievement commands used in
% sections/achievements.tex.
%
% HOW IT WORKS:
% Each \newcommand defines one achievement bullet.
% \Achievement{} wraps the text in an \item so it renders
% as a bullet point inside \begin{ResumeItemize}.
%
% To add a new achievement:
%   1. Define a \newcommand{\YourAchievement} below
%   2. Add \YourAchievement inside the relevant \IfRole{}
%      blocks in sections/achievements.tex
%
% IMPORTANT: Always escape & as \& inside command definitions.
% Raw & is a LaTeX table column separator and will cause a
% "misplaced alignment tab" compile error if unescaped.
% ============================================================


% ============================================================
% COMPETITIVE PROGRAMMING
% ============================================================

% ------------------------------------------------------------
% \LeetCode
% Replace 700+ with your actual solved count.
% \& is used instead of & to avoid a compile error.
% ------------------------------------------------------------
\newcommand{\LeetCode}{
  \Achievement{
    Solved 700+ Data Structures \& Algorithms problems
    across LeetCode, Codeforces, and CodeChef.
  }
}

% ------------------------------------------------------------
% \Codeforces
% NOW wired into \IfSDE in sections/achievements.tex.
% Replace XXXX below with your real max rating before
% compiling. Codeforces ranks for reference:
%   Newbie, Pupil, Specialist, Expert,
%   Candidate Master, Master, International Master, Grandmaster
% ------------------------------------------------------------
\newcommand{\Codeforces}{
  \Achievement{
    Codeforces Specialist (Max Rating: XXXX).
  }
}

% ------------------------------------------------------------
% \CodeChef
% NOT currently called from sections/achievements.tex -- same
% status as \Codeforces above. Fill in your real rating and
% star count, then wire it in when ready.
% ------------------------------------------------------------
\newcommand{\CodeChef}{
  \Achievement{
    CodeChef ★★★ (Max Rating: XXXX).
  }
}


% ============================================================
% ACADEMICS
% ============================================================

% ------------------------------------------------------------
% \CGPAAchievement
% Replace 9.XX with your actual CGPA (e.g. 9.12).
% Used for Quant, Finance, and Consulting roles where
% academic performance is a strong signal.
% ------------------------------------------------------------
\newcommand{\CGPAAchievement}{
  \Achievement{
    Maintained a CGPA of 7.55/10 at VJTI.
  }
}


% ============================================================
% RESEARCH
% ============================================================

% ------------------------------------------------------------
% \ResearchPaper
% Update this if your paper gets published — add the
% conference/journal name and year for stronger credibility.
% ------------------------------------------------------------
\newcommand{\ResearchPaper}{
  \Achievement{
    Developed a novel PCANet-based crop change detection
    methodology using multi-temporal satellite imagery.
  }
}


% ============================================================
% HACKATHONS
% ============================================================

% ------------------------------------------------------------
% \Hackathon
% NOT currently called from sections/achievements.tex --
% "participated, no placement yet" reads as filler in an
% Achievements section, so it's dormant for now. If a specific
% hackathon produced something worth showing (a real build, a
% novel approach, a track win), that's stronger as a Projects
% entry than a line here. Revive this once you have a genuine
% placement (e.g. "Finalist at Smart India Hackathon 2024").
% ------------------------------------------------------------
\newcommand{\Hackathon}{
  \Achievement{
    Participated in national-level hackathons and developed
    end-to-end software solutions under time constraints.
  }
}


% ============================================================
% CERTIFICATIONS
% ============================================================

% ------------------------------------------------------------
% \AWS
% Replace with your actual certification name and year.
% e.g. "AWS Certified Cloud Practitioner (2024)"
% ------------------------------------------------------------
\newcommand{\AWS}{
  \Achievement{
    AWS Certified Cloud Practitioner (2024).
  }
}

% ------------------------------------------------------------
% \GoogleML
% Replace with your actual course/specialization name.
% e.g. "Google Machine Learning Crash Course (2024)"
% ------------------------------------------------------------
\newcommand{\GoogleML}{
  \Achievement{
    Completed Google Machine Learning Crash Course (2024).
  }
}


% ============================================================
% LEADERSHIP
% This generic one-line achievement is retired now that real
% leadership content lives in sections/leadership.tex (Pratibimb
% + Community of Coders). Kept here, commented out, only as a
% record of the old placeholder -- do not re-enable it.
% ============================================================

% \newcommand{\Leadership}{
%   \Achievement{
%     Led student teams in academic and technical projects
%     at VJTI.
%   }
% }


% ============================================================
% OPEN SOURCE
% ============================================================

% ------------------------------------------------------------
% \OpenSource
% NOT currently called from sections/achievements.tex -- no
% real open-source project yet. Fill in the repo/project name
% below once you have one, then wire it into the relevant
% \IfRole{} blocks in sections/achievements.tex.
% ------------------------------------------------------------
\newcommand{\OpenSource}{
  \Achievement{
    Contributed to open-source projects on GitHub.
  }
}


% ============================================================
% GENERAL
% ============================================================

% ------------------------------------------------------------
% \FastLearner
% Currently unused in sections/achievements.tex.
% Add it to a \IfRole{} block there if needed.
% ------------------------------------------------------------
\newcommand{\FastLearner}{
  \Achievement{
    Quickly adapted to new technologies and production
    software development during internship at Wizzer.
  }
}
